\section{Introduction}\label{sec:intro}

An AI coding agent spends nearly all of its money on calls to a large \term{host model}. Between those calls sit many
small choices: which model runs this session, may a file be read without asking the host, how much should the model
reason, is an action safe. A decision model answers such typed questions with a probability per allowed answer in
tens to hundreds of milliseconds. We asked where it earns its place in a harness, and where something simpler does as
well. Routing single queries to cheaper models is well studied~\cite{routellm,frugalgpt}; an agent session is many
calls over one cached, growing context, so we measured whole sessions in pairs and tested the decision model against a
deterministic rule. The report is organised around five takeaways:
\begin{enumerate}[leftmargin=1.6em]
  \item Jev is the best price--quality decision model we tested; OpenAI's Decisions API is the challenger
    (\cref{sec:choose}; preregistered and post hoc).
  \item Route to a cheaper model only when the host is much more expensive once caching and request volume are counted
    (\cref{sec:routing}; preregistered, two hosts).
  \item Configure model and effort once per session; change only after the cache has expired (\cref{sec:once}; once
    preregistered, expiry exploratory and for effort only).
  \item Guard risky actions with a deterministic rule plus the model's plain-text policy; editability is untested
    (\cref{sec:gate}).
  \item Real telemetry measures latency and usage, not accuracy, and showed routing applies far less often than in the
    lab (\cref{sec:telemetry}; exploratory, one owner).
\end{enumerate}
\Cref{sec:setting} describes setting, data and measures once; \cref{sec:recs} gives the configuration and
\cref{sec:conclusion} ranks the takeaways. Every supporting result is in the supplement.
